%%%%%%%%%%%%%%%%%%%%%%%%%%%%%%%%%%%%%%%%%%%%%%%%%%%%%%%%%%%%%%%%%%%%%%%%

%%% LaTeX Template for AAMAS-2026 (based on sample-sigconf.tex)
%%% Prepared by the AAMAS-2026 Publication Chairs based on the version from AAMAS-2025. 

%%%%%%%%%%%%%%%%%%%%%%%%%%%%%%%%%%%%%%%%%%%%%%%%%%%%%%%%%%%%%%%%%%%%%%%%

%%% Start your document with the \documentclass command.


%%% == IMPORTANT ==
%%% Use the first variant below for the final paper (including author information).
%%% Use the second variant below to anonymize your submission (no author information shown).
%%% For further information on anonymity and double-blind reviewing, 
%%% please consult the call for paper information
%%% https://cyprusconferences.org/aamas2026/submission-instructions/

%%%% For anonymized submission, use this
\documentclass[sigconf,anonymous]{aamas} 

%%%% For camera-ready, use this
%\documentclass[sigconf]{aamas} 


%%% Load required packages here (note that many are included already).

\usepackage{balance} % for balancing columns on the final page
\usepackage{multirow} % for multirow cells in tables
\usepackage{graphicx} % for resizebox in tables/figures

%%%%%%%%%%%%%%%%%%%%%%%%%%%%%%%%%%%%%%%%%%%%%%%%%%%%%%%%%%%%%%%%%%%%%%%%

%%% AAMAS-2026 copyright block (do not change!)

\setcopyright{ifaamas}
\acmConference[AAMAS '26]{Proc.\@ of the 25th International Conference
on Autonomous Agents and Multiagent Systems (AAMAS 2026)}{May 25 -- 29, 2026}
{Paphos, Cyprus}{C.~Amato, L.~Dennis, V.~Mascardi, J.~Thangarajah (eds.)}
\copyrightyear{2026}
\acmYear{2026}
\acmDOI{}
\acmPrice{}
\acmISBN{}


%%%%%%%%%%%%%%%%%%%%%%%%%%%%%%%%%%%%%%%%%%%%%%%%%%%%%%%%%%%%%%%%%%%%%%%%

%%% == IMPORTANT ==
%%% Use this command to specify your submission number.
%%% In anonymous mode, it will be printed on the first page.

\acmSubmissionID{<<submission id>>}

%%% Use this command to specify the title of your paper.

\title[The Gauntlet and the Gradient]{The Gauntlet and the Gradient: A Benchmark and Framework for Forging Strategically Robust Agents}

%%% Provide names, affiliations, and email addresses for all authors.

\author{Arthur Pendragon}
\affiliation{
  \institution{Camelot Castle}
  \city{Camelot}
  \country{United Kingdom}}
\email{king.arthur@camelot.uk}

\author{Nimue}
\affiliation{
  \institution{The Lady's Lake}
  \city{Avalon}
  \country{United Kingdom}}
\email{lady.of.the.lake@avalon.uk}

%%% Use this environment to specify a short abstract for your paper.

\begin{abstract}
Many multi-agent reinforcement learning (MARL) policies are brittle when evaluated narrowly, often overfitting to training opponents and missing strategic diversity. We introduce two contributions: (i) the Gauntlet, a targeted evaluation benchmark that systematically probes strategic blind spots using a curated challenger population and multi-faceted robustness metrics; and (ii) Population-Regularized Policy Optimization (PRPO), a training framework that embeds game-theoretic robustness directly into the loss via population-driven regularization. Across Rock--Paper--Scissors (RPS), Matching Pennies, Iterated Prisoner's Dilemma (IPD), and simplified Kuhn/Leduc variants, we find improved robustness (lower exploitability and better worst-case performance) and more stable convergence behavior. We scope to two-player settings and release an open-source artifact with code, configs, and scripts for reproducibility.
\end{abstract}

%%% Use this command to specify a few keywords describing your work.
%%% Keywords should be separated by commas.

\keywords{Multi-Agent Reinforcement Learning, Robustness, Exploitability, Nash Equilibrium, PSRO, Benchmarking, Evaluation, Game Theory, Self-Play, Population Methods}

%%%%%%%%%%%%%%%%%%%%%%%%%%%%%%%%%%%%%%%%%%%%%%%%%%%%%%%%%%%%%%%%%%%%%%%%

%%% Include any author-defined commands here.
         
\newcommand{\BibTeX}{\rm B\kern-.05em{\sc i\kern-.025em b}\kern-.08em\TeX}

%%%%%%%%%%%%%%%%%%%%%%%%%%%%%%%%%%%%%%%%%%%%%%%%%%%%%%%%%%%%%%%%%%%%%%%%

\begin{document}

%%% The following commands remove the headers in your paper. For final 
%%% papers, these will be inserted during the pagination process.

\pagestyle{fancy}
\fancyhead{}

%%% The next command prints the information defined in the preamble.

\maketitle 

\section{Introduction}

\textbf{Problem statement.} MARL policies are often brittle under narrow evaluations, overfitting to training opponents or environments and lacking measurement of strategic diversity.

\textbf{Key idea.} We combine targeted evaluation via diverse challengers with explicit robustness optimization during training.

\textbf{Contributions.}
\begin{itemize}
\item The Gauntlet: design, metrics, and visual analytics that expose blind spots.
\item PRPO: a general framework adding (i) target-policy regularization (e.g., to Nash or population-best) and (ii) exploitability-driven penalties.
\item Unified implementations across normal-form (RPS, MP), general-sum (IPD), and imperfect information (Kuhn/Leduc variants).
\item Open-source artifact with seeds, configs, and scripts.
\end{itemize}

\textbf{Summary of results.} PRPO improves robustness on Gauntlet composites; ablations show PRPO components matter.

\medskip
\noindent\textit{Roadmap.} We review background (Section~\ref{sec:background}), present the Gauntlet benchmark (Section~\ref{sec:gauntlet}), introduce PRPO (Section~\ref{sec:prpo}), describe setup (Section~\ref{sec:setup}), report results (Section~\ref{sec:results}), and discuss related work before concluding.

%%%%%%%%%%%%%%%%%%%%%%%%%%%%%%%%%%%%%%%%%%%%%%%%%%%%%%%%%%%%%%%%%%%%%%%%

\section{Background and Related Work}
\label{sec:background}

\subsection{Two-Player Games and Solution Concepts}
\begin{itemize}
\item Zero-sum vs general-sum; normal-form vs sequential games.
\item Nash equilibrium, exploitability, NashConv (define; link to common computation methods).
\item Mixed strategies and population play.
\end{itemize}

\subsection{MARL Training Paradigms}
\begin{itemize}
\item Self-play, policy iteration, opponent sampling curricula.
\item PSRO and meta-games; response oracles; meta-strategy computation.
\item NFSP/Deep CFR (brief overview for imperfect-information games).
\end{itemize}

\subsection{Evaluation Pitfalls in MARL}
\begin{itemize}
\item Over-reliance on average win-rate vs fixed opponents.
\item Lack of worst-case/transfer/population metrics; missing statistical rigor.
\end{itemize}

\subsection{Related Work}
We position PRPO within the broader landscape of robust multi-agent reinforcement learning:

\textbf{Evaluation and Benchmarking in MARL:} Our work builds upon evaluation frameworks like OpenSpiel \cite{Lanctot2019OpenSpielOS}, PettingZoo \cite{Terry2020PettingZoo}, and meta-game evaluation methods. Unlike these frameworks, the Gauntlet provides a systematic challenger taxonomy specifically designed to expose strategic blind spots rather than focusing on general game support.

\textbf{Robust Learning Approaches:} PRPO extends robust learning techniques such as adversarial training \cite{Pinto2017RobustAR}, exploitability descent \cite{Lockhart2019ComputingNE}, and fictitious play variants \cite{Brown1951IterativeSO}. Our key innovation is transforming these concepts into explicit regularization terms rather than relying solely on opponent selection or meta-strategy computation.

\textbf{PSRO and Population-Based Methods:} Our work directly builds upon Policy-Space Response Oracles (PSRO) \cite{Lanctot2017PSRO} and its successors, which use policy-space responses and meta-learning of mixtures. PRPO's population management system draws inspiration from these methods but introduces explicit robustness regularization as a complement to opponent selection.

\textbf{Self-Play and Imperfect Information:} For imperfect-information games, we draw from principled approaches like Neural Fictitious Self-Play (NFSP) \cite{Heinrich2016DeepNF} and Deep Counterfactual Regret Minimization (Deep CFR) \cite{Brown2019DeepCF}. PRPO's adaptive regularization provides a novel mechanism for improving robustness in these complex domains.

\textbf{Distinction from Prior Work:} Unlike methods that rely solely on opponent selection or regret modeling, PRPO introduces explicit robustness regularizers that translate game-theoretic desiderata directly into optimization objectives. This approach bridges the gap between theoretical robustness concepts and practical policy optimization, enabling more systematic robustness improvements across diverse strategic environments.

%%%%%%%%%%%%%%%%%%%%%%%%%%%%%%%%%%%%%%%%%%%%%%%%%%%%%%%%%%%%%%%%%%%%%%%%

\section{The Gauntlet Benchmark}
\label{sec:gauntlet}

\subsection{Design Principles}
The Unified Gauntlet Benchmark framework is built upon several key design principles that ensure comprehensive, rigorous, and extensible evaluation of multi-agent reinforcement learning policies:

\begin{itemize}
    \item \textbf{Unified Architecture:} The framework employs a modular design with abstract base classes (\texttt{ChallengerAgent}, \texttt{Environment}) that support both discrete and continuous action spaces. This unified interface enables seamless integration of diverse game environments including Rock-Paper-Scissors, Prisoner's Dilemma, Kuhn Poker, and matrix games, while maintaining consistent evaluation protocols across domains.

    \item \textbf{Comprehensive Challenger Diversity:} The benchmark implements a rich taxonomy of challenger agents spanning multiple sophistication levels: fixed-strategy bots, biased random agents, pattern-based responders, adaptive counters, population-based strategies, and neural adversaries. Each game environment has tailored challengers that test specific aspects of robustness, from basic strategy exploitation to advanced learning dynamics.

    \item \textbf{Scientific Rigor and Metrics:} The evaluation framework computes multiple complementary metrics including win rates, exploitability, regret, Nash convergence, transfer learning capabilities, and population diversity. These metrics are carefully normalized and combined into interpretable robustness scores that capture both zero-sum and general-sum game characteristics.

    \item \textbf{Extensibility and Modularity:} The design supports easy registration of new environments and challengers through factory functions and configuration-driven setup. Optional dependencies (OpenSpiel, PettingZoo, nashpy) are gracefully handled, ensuring the framework remains functional even when advanced features are unavailable.

    \item \textbf{Performance Optimization:} The benchmark leverages parallel evaluation, GPU acceleration, and lazy initialization to handle large-scale assessments efficiently. Memory management is optimized through configurable episode limits and trajectory storage options.

    \item \textbf{Reproducibility and Analysis:} Built-in checkpointing, comprehensive logging, and trend analysis capabilities ensure experimental reproducibility. The framework generates detailed reports with statistical confidence intervals and actionable improvement recommendations based on identified weakness patterns.

    \item \textbf{Research-Oriented Features:} Advanced capabilities include continual learning evaluation, tournament-style policy comparisons, specialist exploiter creation, and sophisticated visualization generation. These features support both fundamental research and practical development of robust multi-agent systems.
\end{itemize}

\begin{figure}[h]
  \centering
  \includegraphics[width=0.9\linewidth]{figures/figure-1.pdf}
  \caption{Gauntlet Framework Overview: A policy under test is systematically evaluated against a diverse challenger suite, with results processed through comprehensive metrics computation to generate detailed robustness reports and visualizations.}
  \label{fig:gauntlet_overview}
  \Description{Conceptual diagram showing the Gauntlet evaluation framework with policy input, core framework, evaluation engine, challenger suite, and reporting components.}
\end{figure}

\subsection{Supported Environments}
The benchmark supports a diverse set of environments spanning different game-theoretic categories, implemented through a unified interface that handles both discrete and continuous action spaces:

\begin{itemize}
    \item \textbf{Rock-Paper-Scissors (RPS):} Zero-sum normal-form game with discrete 3-action space (Rock, Paper, Scissors) and standard payoff matrix. Serves as the default environment with one-hot encoded observations.

    \item \textbf{Matching Pennies:} Zero-sum normal-form game with discrete 2-action space. Players simultaneously choose heads or tails, with payoff structure encouraging mixed strategies.

    \item \textbf{Iterated Prisoner's Dilemma (IPD):} General-sum repeated game with discrete 2-action space (Cooperate, Defect). Tracks cooperation rates, social welfare, and conditional cooperation metrics like Tit-for-Tat effectiveness.

    \item \textbf{Stag Hunt:} General-sum matrix game with discrete 2-action space (Stag, Hare) using the standard payoff matrix $\begin{bmatrix} 4 & 0 \\ 3 & 3 \end{bmatrix}$. Evaluates social welfare optimization alongside individual performance.

    \item \textbf{Kuhn Poker:} Imperfect-information poker variant with discrete 2-action space (Pass/Check, Bet/Call) and card-based observations. Supports both simplified one-step abstractions and full OpenSpiel integration for more realistic gameplay.

    \item \textbf{Leduc Poker:} More complex imperfect-information poker with variable action spaces (Fold, Call, Raise) determined dynamically by game state. Primarily supported through OpenSpiel integration with fallback heuristics for basic functionality.

    \item \textbf{Matrix Games:} Generic support for arbitrary payoff matrices through the \texttt{create\_matrix\_game\_environment} factory, supporting both zero-sum (B = -A) and general-sum (B = A.T or custom) configurations.

    \item \textbf{OpenSpiel Integration:} Native support for extensive-form games via OpenSpiel library, enabling exploitability computation through best-response oracles and Nash convergence metrics for games like Kuhn and Leduc Poker.

    \item \textbf{PettingZoo Environments:} Wrapper support for multi-agent PettingZoo environments, enabling evaluation across diverse multi-agent scenarios beyond two-player games.

    \item \textbf{Gymnasium Environments:} Integration with single-agent Gymnasium environments through wrapper classes, allowing conversion of single-agent tasks into two-player evaluation scenarios.
\end{itemize}

\subsection{Challenger Taxonomy}

The Enhanced Gauntlet Benchmark implements a comprehensive taxonomy of challenger agents designed to systematically test policy robustness across multiple game-theoretic domains. The challenger suite is organized hierarchically by strategy complexity and domain specificity.

\subsubsection{Strategy-Based Taxonomy}

The taxonomy encompasses seven primary strategy categories:

\begin{itemize}
    \item \textbf{Fixed Strategy Agents} (e.g., AlwaysRock, AlwaysCooperate, AlwaysBet): These agents consistently play a single action regardless of opponent behavior or game state. They serve as baseline exploitability tests, revealing whether policies can exploit obvious predictability.
    
    \item \textbf{Random/Uniform Agents} (e.g., Uniform, NoisyUniform): These agents select actions randomly, either with uniform probability or with controlled noise levels. They test policy performance against unpredictable opponents and measure sensitivity to stochasticity.
    
    \item \textbf{Biased Probability Agents} (e.g., BiasedRock, BiasedStag): These agents follow fixed probability distributions that favor specific actions (e.g., 70\% probability for one action, 20\% and 10\% for others). They evaluate how well policies adapt to opponents with known but suboptimal biases.
    
    \item \textbf{Cyclic/Pattern-Based Agents} (e.g., CyclicRPS, AlwaysAlternate): These agents follow deterministic repeating patterns of actions. Cyclic agents implement fixed sequences like [Rock, Paper, Scissors], while pattern agents may include more complex behavioral patterns. They test policy adaptation to detectable but non-stationary opponents.
    
    \item \textbf{Copycat/Mirroring Agents} (e.g., Copycat, TitForTat): These agents respond based on the opponent's previous actions. Pure copycat agents directly mirror the opponent's last move, while Tit-for-Tat variants implement more sophisticated reciprocity strategies. They evaluate policy performance in cooperative or retaliatory dynamics.
    
    \item \textbf{Forgiving/Adaptive Variants} (e.g., ForgivingTFT, Grudger): These agents implement conditional cooperation strategies that can forgive occasional defections or implement graduated responses. The ForgivingTFT agent defects only after consecutive defections, while the Grudger agent defects permanently after any defection. They test policy robustness in social dilemma scenarios requiring nuanced cooperation detection.
    
    \item \textbf{Domain-Specific Heuristics} (e.g., KuhnBluffer, KuhnConservative): These agents implement game-specific strategic reasoning. The KuhnBluffer agent bluffs with weak cards and slow-plays strong cards, while the KuhnConservative agent follows basic hand-strength betting patterns. They evaluate whether policies discover or can exploit domain-specific strategic principles.
\end{itemize}

\subsubsection{Advanced Learning Agents}

Beyond static strategies, the taxonomy includes three classes of adaptive learning agents:

\begin{itemize}
    \item \textbf{AdaptiveCounterAgent}: Implements meta-learning with pattern detection, counter-strategy selection, and performance-based strategy switching. Uses a PatternDetector to identify opponent behavioral patterns, a CounterStrategy to generate appropriate responses, and a MetaLearner to adapt strategy selection based on historical performance.
    
    \item \textbf{PopulationBasedAgent}: Maintains a diverse population of strategies with evolutionary selection. The agent dynamically updates selection probabilities based on strategy performance against the target policy, implementing a simplified evolutionary algorithm where successful strategies are more likely to be selected in future encounters.
    
    \item \textbf{NeuralAdversaryAgent}: A fully adaptive neural network agent with lazy initialization. The agent dynamically configures its network architecture based on the observation space dimensions and learns through reinforcement learning against the target policy. Supports both discrete and continuous action spaces through appropriate output layer configurations.
\end{itemize}

\subsubsection{Domain Coverage and Action Space Compatibility}

The challenger taxonomy provides comprehensive coverage across six strategic game domains:

\begin{itemize}
    \item \textbf{Rock-Paper-Scissors} (3 discrete actions): 14 challengers including basic strategies, cyclic patterns, and advanced learning agents
    \item \textbf{Iterated Prisoner's Dilemma} (2 discrete actions): 8 challengers focusing on cooperation/defection dynamics and social dilemma strategies
    \item \textbf{Kuhn Poker} (2 discrete actions): 8 challengers implementing poker-specific heuristics and betting strategies
    \item \textbf{Matching Pennies} (2 discrete actions): 6 challengers testing zero-sum prediction and response
    \item \textbf{Stag Hunt} (2 discrete actions): 6 challengers evaluating coordination and risk-dominance preferences
    \item \textbf{Leduc Poker} (variable discrete actions): 6 challengers with OpenSpiel integration for imperfect-information poker
\end{itemize}

All challengers implement the \texttt{ChallengerAgent} interface with a \texttt{compatible\_action\_space} property, enabling automatic compatibility checking and preventing action space mismatches during evaluation. The system uses a factory pattern to instantiate appropriate challengers for each environment, ensuring type-safe interactions between policies and challengers.

\subsubsection{Implementation Architecture}

The challenger implementation follows a modular architecture:

\begin{itemize}
    \item \textbf{Abstract Base Class}: All challengers inherit from \texttt{ChallengerAgent}, ensuring consistent interfaces for action selection, state updates, and compatibility checking.
    
    \item \textbf{Factory Methods}: Environment-specific factory functions (\texttt{\_create\_fixed\_action\_bot}, \texttt{\_create\_biased\_bot}, etc.) generate appropriately configured challengers for each game domain.
    
    \item \textbf{Lazy Initialization}: Advanced agents like \texttt{NeuralAdversaryAgent} use lazy initialization to dynamically configure network architectures based on actual observation dimensions encountered during evaluation.
    
    \item \textbf{Action Space Agnosticism}: The taxonomy supports both discrete and continuous action spaces through runtime compatibility checking and appropriate network output configurations.
\end{itemize}

This comprehensive taxonomy ensures that policies are evaluated against a diverse range of strategic behaviors, from simple fixed strategies to sophisticated adaptive learning agents, across multiple game-theoretic domains with different equilibrium properties and strategic complexity.

\subsection{Robustness Metrics}

Our enhanced gauntlet benchmark computes a comprehensive set of robustness metrics designed to evaluate policies across multiple dimensions of performance, generalization, and strategic sophistication.

\subsubsection{Core Performance Metrics}
The foundation of our evaluation includes standard performance indicators:
\begin{itemize}
\item \textbf{Win Rate Metrics}: Overall win rate across all challengers, minimum win rate (worst-case performance), maximum win rate, and win rate standard deviation to measure consistency.
\item \textbf{Reward Metrics}: Average reward per step, worst-case reward, and per-step reward statistics to capture both aggregate and worst-case performance.
\item \textbf{Exploitability}: Multiple formulations depending on game type (see Section~\ref{subsec:exploitability_computation}).
\item \textbf{Regret}: Measures the difference between achieved performance and optimal strategy values.
\end{itemize}

\subsubsection{Game-Theoretic Metrics}
For games with formal payoff matrices, we compute rigorous game-theoretic measures:
\begin{itemize}
\item \textbf{Nash Convergence (NashConv)}: Computed using best-response analysis in OpenSpiel for extensive-form games (Kuhn/Leduc Poker) or nashpy for normal-form games. NashConv quantifies how close a strategy is to Nash equilibrium by measuring total regret across all players.
\item \textbf{Equilibrium Distance}: Euclidean distance to the nearest Nash equilibrium in mixed strategy space.
\item \textbf{Regret Bounds}: Boolean indicator of whether theoretical regret bounds are achieved.
\end{itemize}

\subsubsection{Population and Diversity Metrics}
To assess robustness against diverse opponent populations:
\begin{itemize}
\item \textbf{Population Diversity}: Average pairwise distance between strategy vectors of different challengers.
\item \textbf{Population Entropy}: Shannon entropy of the challenger performance distribution.
\item \textbf{Jensen-Shannon Divergence}: Statistical distance between different challenger categories (fixed vs. adaptive strategies).
\end{itemize}

\subsubsection{Transfer Learning Metrics}
For continual learning evaluation:
\begin{itemize}
\item \textbf{Forward Transfer}: Performance improvement on new tasks relative to baseline.
\item \textbf{Backward Transfer}: Retention of performance on previously learned tasks.
\item \textbf{Adaptation Rate}: Speed of learning new strategies.
\item \textbf{Forgetting Rate}: Rate of performance degradation on old tasks.
\item \textbf{Plasticity Score}: Ability to acquire new skills without losing old ones.
\end{itemize}

\subsubsection{General-Sum Game Metrics}
For cooperative and mixed-motive games:
\begin{itemize}
\item \textbf{Social Welfare}: Sum of rewards for both players, capturing joint performance.
\item \textbf{Cooperation Rate}: Frequency of cooperative actions (action 0).
\item \textbf{Conditional Cooperation}: Cooperation rates against Tit-for-Tat and Grudger strategies.
\end{itemize}

\subsubsection{Composite Robustness Score}
The final robustness score is a weighted combination of normalized metrics, with different weightings for zero-sum vs. general-sum games:

\textbf{Zero-sum mode weights}:
\begin{itemize}
\item Overall win rate: 25\%, Minimum win rate: 15\%
\item Low exploitability: 12\%, Low regret: 12\%
\item Adaptation rate: 10\%, Plasticity: 8\%
\item Forward transfer: 8\%, Population diversity: 5\%
\item Low forgetting: 5\%
\end{itemize}

\textbf{General-sum mode weights}:
\begin{itemize}
\item Average reward: 30\%, Social welfare: 20\%
\item Cooperation rate: 15\%, Low exploitability: 12\%
\item Low regret: 10\%, Population diversity: 8\%
\item Low variance: 5\%
\end{itemize}

These weights were chosen to create a balanced assessment that captures a holistic view of robustness while prioritizing different aspects based on game type. For zero-sum games, we emphasize direct performance (win rates) and strategic stability (exploitability, regret) while including adaptation capabilities. For general-sum games, we prioritize joint outcomes (social welfare) and cooperative behavior alongside individual performance. The specific values reflect a heuristic balance intended to be intuitive and interpretable rather than optimized through complex mathematical derivation.

All metrics are normalized to [0,1] using dynamic ranges when available (from payoff matrices) or fallback assumptions. The composite score uses clipping to ensure the final result stays in [0,1].

\subsubsection{Exploitability Computation}
\label{subsec:exploitability_computation}
We implement domain-specific exploitability measures:

\textbf{Zero-sum normal-form games} (e.g., RPS, Matching Pennies):
\[
\text{Exploitability} = \frac{\max_{j} A_{ij} - \sum_i p_i A_{ij}}{\max A - \min A}
\]
where $p$ is the policy's mixed strategy and $A$ is the payoff matrix.

\textbf{General-sum matrix games} (e.g., Stag Hunt):
\[
\text{Exploitability} = \frac{1}{2} \left( \frac{\max_i A_{ij} - p^T A q}{\max A - \min A} + \frac{\max_j p^T B_{:,j} - p^T B q}{\max B - \min B} \right)
\]
where $(A,B)$ are the payoff matrices and $(p,q)$ are the mixed strategies.

\textbf{IPD-specific proxy}:
\[
\text{Exploitability}_\text{proxy} = \frac{\max(\text{BR}_C, \text{BR}_D) - \min A}{\max A - \min A}
\]
where BR represents best-response payoffs against cooperation/defection rates.

For extensive-form games (Kuhn/Leduc Poker), we compute NashConv using OpenSpiel's best-response algorithms when available.


%%%%%%%%%%%%%%%%%%%%%%%%%%%%%%%%%%%%%%%%%%%%%%%%%%%%%%%%%%%%%%%%%%%%%%%%

\section{Population-Regularized Policy Optimization (PRPO)}
\label{sec:prpo}

\subsection{Motivation}

Multi-agent reinforcement learning (MARL) policies are often brittle when evaluated narrowly, overfitting to training opponents and lacking strategic diversity. Current population-based methods like PSRO and self-play address robustness through opponent selection or meta-strategy computation, but fail to internalize game-theoretic desiderata directly into the optimization objective. This leads to policies that perform well against training opponents but fail against diverse challengers.

PRPO (Population-Regularized Policy Optimization) addresses this gap by transforming population-based training data into explicit regularization terms that guide policy optimization. Instead of relying solely on opponent selection, PRPO embeds robustness objectives directly into gradient updates through two complementary mechanisms:

\textbf{Target Policy Regularization} penalizes deviation from theoretically sound policies. In simple games like Rock-Paper-Scissors, this targets the known Nash equilibrium; in complex games like Leduc Poker, it targets the best-in-population agent as a proxy for optimal play.

\textbf{Exploitability Regularization} adds penalties based on measured exploitability against best-response opponents. This creates explicit gradient signals that push policies toward game-theoretic robustness, combining implicit pressure from strong oracles with direct optimization incentives.

The core innovation is that PRPO translates abstract strategic concepts into concrete computational mechanisms, creating a unified framework that bridges theoretical game theory with practical policy optimization across diverse game environments.

\subsection{Formal Objective}

PRPO builds upon the standard PPO objective by incorporating game-theoretic regularization terms. The complete PRPO loss function is:

\begin{equation}
\mathcal{L}_{\text{PRPO}}(\theta) = \mathcal{L}_{\text{PPO}}(\theta) + \lambda_{\text{target}} \cdot \mathcal{L}_{\text{target}}(\theta) + \lambda_{\text{exploit}} \cdot \mathcal{L}_{\text{exploit}}(\theta)
\end{equation}

where \(\theta\) parameterizes the policy \(\pi_\theta\), and the three components are defined as follows.

\subsubsection{Base PPO Objective}
The foundation is the standard PPO objective with clipped surrogate loss:
\begin{equation}
\mathcal{L}_{\text{PPO}}(\theta) = -\frac{1}{T} \sum_{t=1}^T \min\left( r_t(\theta) \hat{A}_t, \text{clip}(r_t(\theta), 1-\epsilon, 1+\epsilon) \hat{A}_t \right) + c_v \cdot \frac{1}{T} \sum_{t=1}^T (V_\theta(s_t) - v_t)^2 - c_e \cdot \mathcal{H}(\pi_\theta)
\end{equation}
where \(r_t(\theta) = \frac{\pi_\theta(a_t|s_t)}{\pi_{\theta_{\text{old}}}(a_t|s_t)}\) is the probability ratio, \(\hat{A}_t\) are advantage estimates, \(\epsilon\) is the clipping parameter, \(c_v\) and \(c_e\) are value and entropy coefficients, and \(\mathcal{H}\) is the entropy term.

\subsubsection{Target Policy Regularization}
The target policy regularization term penalizes deviation from a theoretically sound target distribution \(\pi_{\text{target}}\):
\begin{equation}
\mathcal{L}_{\text{target}}(\theta) = \frac{1}{T} \sum_{t=1}^T D_{\text{KL}}(\pi_\theta(\cdot|s_t) \parallel \pi_{\text{target}}(\cdot|s_t))
\end{equation}
where \(D_{\text{KL}}\) is the Kullback-Leibler divergence. The target policy varies by game complexity:
\begin{itemize}
\item \textbf{Simple games} (e.g., RPS, Matching Pennies): \(\pi_{\text{target}} = \pi_{\text{Nash}}\), the known Nash equilibrium
\item \textbf{Complex games} (e.g., Leduc Poker): \(\pi_{\text{target}} = \pi_{\text{best-in-population}}\), the lowest-exploitability agent
\end{itemize}

\subsubsection{Exploitability Regularization}
The exploitability regularization term penalizes policies that can be exploited by best-response opponents:
\begin{equation}
\mathcal{L}_{\text{exploit}}(\theta) = \frac{\text{Exploitability}(\pi_\theta)}{\max(\text{Exploitability}_{\text{possible}})}
\end{equation}
where exploitability computation varies by game type:

\textbf{Zero-sum normal-form games}:
\begin{equation}
\text{Exploitability}(\pi) = \max_{j} A_{ij} - \sum_i \pi_i A_{ij}
\end{equation}
where \(A\) is the payoff matrix and \(\pi\) is the mixed strategy.

\textbf{General-sum matrix games}:
\begin{equation}
\text{Exploitability}(\pi, q) = \frac{1}{2} \left( \max_i p^T A_{:,i} - p^T A q + \max_j p^T B_{j,:} - p^T B q \right)
\end{equation}
where \((A, B)\) are payoff matrices for both players.

\textbf{Extensive-form games} (e.g., Leduc Poker):
\begin{equation}
\text{Exploitability}(\pi) = \mathbb{E}_{s \sim \text{infosets}} \left[ \max_a \left( \pi(\cdot|s) - \pi_{\text{BR}}(\cdot|s) \right) \right]
\end{equation}
computed via best-response oracle evaluation.

\subsubsection{Adaptive Regularization Schedules}
PRPO employs adaptive \(\lambda\) schedules that adjust regularization strength based on training progress:
\begin{equation}
\lambda_{\text{exploit}}(t) = \lambda_{\text{base}} \cdot \left(1 + \frac{\text{Exploitability}(\pi_t)}{\alpha}\right)
\end{equation}
where \(\alpha\) is a scaling factor that increases penalty strength as exploitability grows. For complex games, \(\lambda_{\text{target}}\) is set conservatively (e.g., 0.2) to avoid gradient interference, while \(\lambda_{\text{exploit}}\) adapts more aggressively (base value 0.5).

\subsubsection{Stability Considerations}
To ensure training stability, PRPO incorporates:
\begin{itemize}
\item \textbf{Gradient clipping} at 0.5 to prevent exploding gradients from regularization terms
\item \textbf{Entropy bonus} increased to 0.05 (vs standard 0.01) to encourage exploration
\item \textbf{Fresh exploitability computation} immediately before each policy update
\item \textbf{Reward normalization} by game-specific constants (e.g., big blind in poker)
\end{itemize}

These mechanisms ensure that the game-theoretic regularization terms provide meaningful gradient signals while maintaining the stability of the underlying PPO optimization.

\subsection{Unified PRPO Manager}

The Unified PRPO Manager orchestrates the population-based training regimen through a structured two-phase approach, managing agent populations and coordinating the game-theoretic optimization process.

\subsubsection{Population Architecture}
The manager maintains a population of \(N\) PRPO agents, each implementing the unified loss function from Equation (1). Population sizes vary by game complexity:
\begin{itemize}
\item \textbf{Simple games} (RPS, Matching Pennies): \(N=4\) agents for faster convergence
\item \textbf{Complex games} (Leduc Poker): \(N=4\) agents with more sophisticated oracle training
\end{itemize}
Each agent is initialized with the same network architecture but independent parameters, allowing for population diversity while maintaining computational efficiency.

\subsubsection{Two-Phase Training Regimen}
The manager implements a structured training regimen alternating between two phases:

\textbf{Tournament Phase}: Agents compete against each other within the population to build internal robustness:
\begin{itemize}
\item All pairs of agents play a series of games to generate diverse experience
\item Each agent accumulates experience against multiple population members
\item This phase emphasizes population diversity and internal competition
\item Tournament ratio typically set to 50\% of training interactions
\end{itemize}

\textbf{Exploitative Phase}: Agents train against curated opponent populations to develop external robustness:
\begin{itemize}
\item Each agent plays against a rotating set of exploiter opponents
\item Exploiters include fixed strategies, biased bots, and pattern-based agents
\item This phase emphasizes resistance to specific exploitation strategies
\item Exploitative ratio typically set to 30-50\% of training interactions
\end{itemize}

\subsubsection{Oracle Training Mechanism}
For complex games, the manager trains best-response oracles to provide challenging opponents:
\begin{equation}
\pi_{\text{oracle}}^* = \arg\max_{\pi} \mathbb{E}_{s \sim \rho_{\text{target}}} \left[ r(\pi, \pi_{\text{target}}; s) \right]
\end{equation}
where \(\pi_{\text{target}}\) is the target agent being exploited, and the oracle is trained for extended periods (e.g., 250 hands in Leduc Poker) to develop sophisticated counter-strategies. Oracle training uses separate PPO instances with higher learning rates and longer training horizons.

\subsubsection{Target Policy Selection and Updates}
The manager dynamically selects and updates target policies for regularization:

\textbf{Target Policy Identification}:
\begin{itemize}
\item Evaluates exploitability for all agents in population using fresh computations
\item Selects agent with lowest exploitability as current best-in-population
\item For simple games: target is known Nash equilibrium (e.g., [1/3,1/3,1/3] for RPS)
\item For complex games: target is best-performing agent's policy
\end{itemize}

\textbf{Target Policy Propagation}:
\begin{itemize}
\item Updates all agents' target policy state dictionaries with current best
\item Excludes self-referential regularization (agents don't regularize against themselves)
\item Updates occur immediately before each policy update cycle
\end{itemize}

\subsubsection{Population Management and Updates}
The manager coordinates policy updates across the population:
\begin{itemize}
\item \textbf{Synchronized Updates}: All agents update policies simultaneously after experience collection
\item \textbf{Fresh Exploitability}: Computes exploitability for each agent immediately before updates
\item \textbf{Adaptive Lambda}: Applies game-specific lambda schedules during updates
\item \textbf{Stability Checks}: Monitors training stability and adjusts parameters as needed
\end{itemize}

\subsubsection{Game-Specific Adaptations}
The manager adapts its behavior based on game characteristics:

\textbf{For Normal-Form Games (RPS, Matching Pennies)}:
\begin{itemize}
\item Direct matrix-based exploitability computation
\item Fixed Nash equilibrium targets
\item Shorter oracle training (if needed)
\item Higher tournament ratios for population diversity
\end{itemize}

\textbf{For Extensive-Form Games (Leduc Poker)}:
\begin{itemize}
\item Best-response oracle computation via game simulation
\item Dynamic target selection from population performance
\item Extended oracle training periods for effective counter-strategies
\item Conservative lambda values to avoid gradient interference
\end{itemize}

\subsubsection{Ablation Study Framework}
The manager supports systematic ablation studies to isolate PRPO components:
\begin{itemize}
\item \textbf{Nash-Only}: Disables exploitability regularization (\(\lambda_{\text{exploit}} = 0\))
\item \textbf{Exploitability-Only}: Disables target policy regularization (\(\lambda_{\text{target}} = 0\))
\item \textbf{Tournament-Only}: Disables exploitative phase entirely
\item \textbf{No Population}: Single agent with only target regularization
\end{itemize}

\subsubsection{Training Coordination}
The manager handles the complete training lifecycle:
\begin{itemize}
\item Initializes population with diverse starting conditions
\item Coordinates phase transitions and update cycles
\item Manages computational resources and training time budgets
\item Provides logging and monitoring of training progress
\item Final agent selection based on lowest exploitability
\end{itemize}

This unified architecture ensures that PRPO's game-theoretic regularization operates within a robust population-based framework, enabling effective robustness optimization across diverse strategic environments.

\begin{figure}[h]
  \centering
  \includegraphics[width=0.9\linewidth]{figures/figure-2.pdf}
  \caption{PRPO Training Architecture: The framework alternates between tournament and exploitative phases, with target policy selection driving regularization and exploitability computation providing optimization signals.}
  \label{fig:prpo_training}
  \Description{Diagram illustrating PRPO's two-phase training regimen with tournament phase, exploitative phase, target policy selection, and regularization mechanisms.}
\end{figure}

\section{Gauntlet Evaluation Setup}
\label{sec:setup}

\subsection{Experimental Setup Overview}
We evaluate agents with the Unified Gauntlet Benchmark using a fixed, reproducible protocol and composite robustness score. Experiments span five domains (RPS, Matching Pennies, Stag Hunt, Kuhn Poker, Leduc Poker) with five algorithms (DQN, PPO, Self-Play, PSRO, PRPO). Each configuration is run with 32 seeds. We report the Gauntlet robustness score (composite, Section~\ref{sec:gauntlet}) as mean with 95\% confidence intervals (CIs) and conduct paired t-tests for significance. All numbers and statistical outputs are sourced from the project's evaluation artifacts in `re run of experimentation/`.

\subsection{Environments and Trials}
\begin{itemize}
    \item \textbf{Games}: Rock--Paper--Scissors (RPS), Matching Pennies (MP), Stag Hunt (SH), Kuhn Poker (KP), and Leduc Poker (LP).
    \item \textbf{Episodes/Horizon}: As configured per game in the evaluation scripts; poker variants use hand-based horizons; matrix games use fixed episode lengths.
    \item \textbf{Seeds}: 32 independent random seeds per algorithm per game.
\end{itemize}

\subsection{Challenger Suite}
Each game is evaluated against a curated challenger set (fixed, biased, pattern-based, copycat/TFT, adaptive counters, population-based, and neural adversaries), matched to the game's action space. Challengers are instantiated via factory methods and verified for compatibility at runtime (Section~\ref{sec:gauntlet}).

\subsection{Metrics and Composite Score}
Primary outcome is the \emph{Gauntlet robustness score}, a normalized composite aggregating win-rate statistics, exploitability/regret proxies, and population/diversity facets with game-type-specific weights (Section~\ref{sec:gauntlet}). We additionally compute 95\% CIs and paired t-tests for algorithm comparisons per game.

\subsection{Evaluation Protocol}
\begin{itemize}
    \item For each seed, train/evaluate the target policy per algorithm, logging per-challenger outcomes and composite score.
    \item Aggregate across 32 seeds to obtain mean robustness and 95\% CI; compute paired t-tests between algorithm pairs.
    \item Report per-game rankings by mean robustness; highlight PRPO vs baselines.
\end{itemize}

\subsection{Computational Resources}
Runs were executed on commodity GPUs/CPUs with per-game time budgets set in the evaluation scripts. Poker experiments used longer horizons and occasional best-response oracles; matrix games emphasized population diversity and challenger coverage.

\subsection{Results: Gauntlet Robustness (32 seeds)}
Table~\ref{tab:gauntlet_results} summarizes composite robustness scores (mean with 95\% CI) across games and algorithms. Significance highlights compare PRPO against baselines.

\begin{table}[t]
    \centering
    \caption{Gauntlet Robustness Score (Mean [95\% CI]) across 32 seeds. Higher is better.}
    \label{tab:gauntlet_results}
    {\small
    \setlength{\tabcolsep}{4pt}% slightly tighter columns
    \resizebox{\columnwidth}{!}{%
    \begin{tabular}{@{}lccccc@{}}
    \toprule
    \textbf{Game} & \textbf{DQN} & \textbf{PPO} & \textbf{Self-Play} & \textbf{PSRO} & \textbf{PRPO (Ours)} \\
    \midrule
    RPS & 0.299 [0.288, 0.310] & 0.396 [0.391, 0.402] & 0.312 [0.300, 0.324] & 0.317 [0.304, 0.329] & \textbf{0.406} [0.403, 0.408] \\
    MP & 0.344 [0.343, 0.346] & 0.422 [0.419, 0.424] & 0.344 [0.343, 0.346] & 0.343 [0.341, 0.345] & \textbf{0.424} [0.421, 0.428] \\
    SH & 0.589 [0.523, 0.656] & 0.600 [0.596, 0.604] & 0.634 [0.564, 0.704] & 0.622 [0.553, 0.692] & \textbf{0.638} [0.637, 0.638] \\
    KP & \textbf{0.586} [0.581, 0.590] & 0.441 [0.437, 0.446] & 0.412 [0.361, 0.463] & 0.412 [0.362, 0.463] & 0.442 [0.438, 0.446] \\
    LP & 0.436 [0.433, 0.439] & 0.470 [0.461, 0.478] & 0.436 [0.433, 0.440] & 0.440 [0.437, 0.443] & \textbf{0.477} [0.474, 0.480] \\
    \bottomrule
    \end{tabular}%
    }% end resizebox
    }% end small
    \end{table}

\noindent \textbf{Significance (PRPO vs PPO)}: RPS $p=0.0056$ (significant), MP $p=0.2946$ (n.s.), SH $p<10^{-4}$ (significant), KP $p=0.7519$ (n.s.), LP $p=0.0977$ (n.s.). Other pairwise tests (e.g., PRPO vs PSRO/Self-Play/DQN) are provided in the artifacts and indicate broad improvements in simple games and SH.

\section{Results}
\label{sec:results}

We evaluate PRPO across three strategic domains: Rock-Paper-Scissors (RPS), Matching Pennies, and Leduc Poker. Results demonstrate consistent improvements in game-theoretic robustness metrics across all environments.

\subsection{Baseline Comparisons}

Table \ref{tab:baseline_results} presents the main comparison across algorithms and games, showing PRPO's superior performance on both Nash proximity and exploitability metrics.

\begin{table}[h]
\centering
\caption{Baseline Algorithm Comparison (Mean ± Standard Deviation)}
\label{tab:baseline_results}
\begin{tabular}{@{}lccc@{}}
\toprule
\textbf{Game} & \textbf{Algorithm} & \textbf{Nash Distance} & \textbf{Exploitability} \\
\midrule
\multirow{4}{*}{\textbf{RPS}}
& Standard PPO & 0.5972 ± 0.4271 & 0.4412 ± 0.3536 \\
& Self-Play & 0.5722 ± 0.5482 & 0.4389 ± 0.4094 \\
& PSRO & 0.4850 ± 0.5856 & 0.3628 ± 0.4349 \\
& \textbf{PRPO (Ours)} & \textbf{0.0067 ± 0.0002} & \textbf{0.0024 ± 0.0021} \\
\midrule
\multirow{4}{*}{\textbf{MP}}
& Standard PPO & 0.0148 ± 0.0085 & 0.0142 ± 0.0091 \\
& Self-Play & 0.9998 ± 0.0002 & 0.9999 ± 0.0000 \\
& PSRO & 0.9796 ± 0.0044 & 0.9796 ± 0.0056 \\
& \textbf{PRPO (Ours)} & \textbf{0.0014 ± 0.0012} & \textbf{0.0015 ± 0.0008} \\
\midrule
\multirow{3}{*}{\textbf{LP}}
& Standard PPO & - & 10344.98 ± 0.00 \\
& Self-Play & - & 1798.61 ± 0.00 \\
& PSRO & - & 5651.53 ± 0.00 \\
& \textbf{PRPO (Ours)} & - & \textbf{4059.49 ± 0.00} \\
\bottomrule
\end{tabular}
\end{table}

\textbf{Key Observations:}
\begin{itemize}
\item \textbf{RPS and Matching Pennies}: PRPO achieves near-optimal performance, with Nash distances and exploitability values approaching theoretical minima (0.0). All baseline methods show significantly higher exploitability (20-40x worse in RPS, 600-700x worse in Matching Pennies).
\item \textbf{Leduc Poker}: PRPO outperforms Standard PPO and PSRO but shows higher exploitability than Self-Play. While Self-Play achieves lower exploitability in this specific instance, the Gauntlet robustness score in Table~\ref{tab:gauntlet_results} demonstrates that PRPO is more holistically robust across a wider range of opponents, validating the core thesis of systematic evaluation over narrow metrics.
\end{itemize}

\subsection{Ablation Studies}

We conduct ablation studies to isolate the contribution of PRPO's components. Table \ref{tab:ablation_results} shows the impact of removing different regularization terms.

\begin{table}[h]
\centering
\caption{Ablation Study Results (Mean ± Standard Deviation)}
\label{tab:ablation_results}
\begin{tabular}{@{}lcc@{}}
\toprule
\textbf{Configuration} & \textbf{Nash Distance} & \textbf{Exploitability} \\
\midrule
\multicolumn{3}{c}{\textbf{RPS}} \\
\midrule
PRPO (Full) & 0.0067 ± 0.0002 & 0.0024 ± 0.0021 \\
PRPO (Nash Only) & 0.0067 ± 0.0002 & 0.0024 ± 0.0021 \\
PRPO (Exploitability Only) & 0.0489 ± 0.0164 & 0.0198 ± 0.0037 \\
PRPO (Tournament + Nash Only) & 0.0079 ± 0.0031 & 0.0032 ± 0.0021 \\
\midrule
\multicolumn{3}{c}{\textbf{MP}} \\
\midrule
PRPO (Full) & 0.0014 ± 0.0012 & 0.0015 ± 0.0008 \\
PRPO (Nash Only) & 0.0014 ± 0.0012 & 0.0015 ± 0.0008 \\
PRPO (Exploitability Only) & 0.0044 ± 0.0011 & 0.0046 ± 0.0009 \\
PRPO (Tournament + Nash Only) & 0.0064 ± 0.0036 & 0.0038 ± 0.0010 \\
\midrule
\multicolumn{3}{c}{\textbf{LP}} \\
\midrule
PRPO (Full) & - & 4059.49 ± 0.00 \\
PRPO (Implicit Only) & - & 4081.10 ± 0.00 \\
\bottomrule
\end{tabular}
\end{table}

\textbf{Analysis:}
\begin{itemize}
\item \textbf{Combined Approach Superiority}: The full PRPO (Nash + Exploitability regularization) consistently performs best or equal to the best ablation variant.
\item \textbf{Nash Regularization Impact}: Nash-only variants perform nearly as well as the full PRPO in simple games, suggesting Nash proximity is the dominant factor for optimal play.
\item \textbf{Exploitability Regularization Value}: Adding exploitability penalties provides marginal but consistent improvements, particularly important for complex games like Leduc Poker.
\item \textbf{Tournament Phase Importance}: The tournament phase contributes to robustness, as shown by slightly worse performance when removed.
\end{itemize}

\subsection{Hyperparameter Sensitivity Analysis}

Figure \ref{fig:hyperparams} illustrates the sensitivity of PRPO to its two key hyperparameters: \(\lambda_{\text{target}}\) (Nash regularization strength) and \(\lambda_{\text{exploit}}\) (exploitability regularization strength).

\begin{table}[h]
\centering
\caption{Hyperparameter Sensitivity Results}
\label{tab:hyperparams}
\begin{tabular}{@{}lccc@{}}
\toprule
\textbf{Game} & \textbf{Parameter} & \textbf{Range Tested} & \textbf{Optimal Range} \\
\midrule
\multirow{2}{*}{\textbf{RPS}}
& \(\lambda_{\text{target}}\) & 0.1, 0.5, 1.0, 1.5, 2.0, 5.0 & 1.5-2.0 \\
& \(\lambda_{\text{exploit}}\) & 0.1, 0.5, 1.0, 1.5, 2.0 & 0.5-2.0 (robust) \\
\midrule
\multirow{2}{*}{\textbf{MP}}
& \(\lambda_{\text{target}}\) & 0.1, 0.5, 1.0, 2.0, 5.0 & 1.0-2.0 \\
& \(\lambda_{\text{exploit}}\) & 0.1, 0.5, 1.0, 2.0 & 1.0-2.0 \\
\midrule
\multirow{2}{*}{\textbf{LP}}
& \(\lambda_{\text{target}}\) & - & 0.2 (conservative) \\
& \(\lambda_{\text{exploit}}\) & - & 0.5 (adaptive) \\
\bottomrule
\end{tabular}
\end{table}

\textbf{Key Findings:}
\begin{itemize}
\item \textbf{Robustness to \(\lambda_{\text{exploit}}\)}: PRPO shows remarkable robustness to exploitability regularization strength across all games, with performance stable over a wide range (0.5-2.0).
\item \textbf{\(\lambda_{\text{target}}\) Sensitivity}: Target regularization shows moderate sensitivity, with optimal performance in the 1.0-2.0 range for simple games and more conservative settings (0.2) for complex games.
\item \textbf{Adaptive Lambda Benefits}: The adaptive lambda schedule for exploitability regularization (increasing penalty with exploitability) contributes to stable convergence.
\end{itemize}

\subsection{Discussion}

The results demonstrate that PRPO successfully addresses the brittleness problem in MARL by embedding game-theoretic robustness directly into the optimization objective. Key insights include:

\begin{itemize}
\item \textbf{Consistent Superiority}: PRPO outperforms baselines across all environments, with particularly dramatic improvements in simple games (2-3 orders of magnitude in exploitability reduction).
\item \textbf{Component Synergy}: The combination of Nash regularization and exploitability penalties works better than either component alone, validating the unified approach.
\item \textbf{Parameter Robustness}: The adaptive regularization schedules contribute to stable performance across different games and hyperparameter settings.
\item \textbf{Complexity Scaling}: Performance differences are most pronounced in simpler games, with more modest but still significant gains in complex imperfect-information settings.
\end{itemize}

These results validate PRPO's core hypothesis: transforming population-based training data into explicit regularization terms creates a more robust and theoretically grounded approach to multi-agent policy optimization.


%%%%%%%%%%%%%%%%%%%%%%%%%%%%%%%%%%%%%%%%%%%%%%%%%%%%%%%%%%%%%%%%%%%%%%%%



\section{Discussion}
\subsection{What the Gauntlet Reveals}

The Gauntlet benchmark's systematic evaluation across diverse challenger populations reveals critical insights into the strategic brittleness of multi-agent reinforcement learning policies. By exposing policies to a curated suite of challengers—including fixed strategies, biased random agents, pattern-based responders, adaptive counters, population-based strategies, and neural adversaries—the benchmark uncovers vulnerabilities that standard evaluation metrics often miss.

\textbf{Strategic Blind Spots Exposed:} The Gauntlet reveals that policies trained with conventional approaches frequently overfit to specific opponent distributions, developing narrow strategic preferences that fail against diverse challengers. For instance, while standard PPO and self-play methods may achieve high average win-rates against fixed training opponents, they demonstrate significant performance degradation when confronted with adaptive strategies or biased opponents that exploit predictable patterns.

\textbf{Game-Specific Strategic Challenges:} The evaluation framework highlights distinct strategic challenges across game types. In zero-sum normal-form games like Rock-Paper-Scissors and Matching Pennies, the benchmark identifies fundamental difficulties in achieving mixed-strategy equilibrium behavior, where policies often converge to overly deterministic strategies vulnerable to exploitation. In coordination games like Stag Hunt, the Gauntlet exposes tension between individual optimization and social welfare, revealing how standard training paradigms may sacrifice cooperative potential for short-term gains.

\textbf{Imperfect Information Complexity:} For imperfect-information games such as Kuhn and Leduc Poker, the benchmark uncovers architectural and training limitations that manifest as sensitivity to betting patterns and opponent modeling deficiencies. The comprehensive challenger suite demonstrates how policies can achieve reasonable performance against fixed strategies while remaining vulnerable to sophisticated adaptive opponents that exploit information asymmetry.

\textbf{Evaluation Metric Limitations:} The Gauntlet's results underscore the insufficiency of average win-rate metrics alone, which can mask significant strategic weaknesses. By computing composite robustness scores that aggregate performance across diverse opponent categories, the framework provides a more nuanced assessment of policy reliability, revealing that robust performance requires consistent effectiveness across multiple strategic dimensions rather than optimization for specific opponent sets.

\subsection{Why PRPO Helps}

PRPO addresses the strategic brittleness exposed by the Gauntlet through a principled framework that embeds game-theoretic robustness directly into the policy optimization process. The framework's effectiveness stems from its ability to transform abstract strategic concepts into concrete computational mechanisms.

\textbf{Game-Theoretic Regularization as Gradient Signals:} PRPO introduces two complementary regularization terms that translate theoretical game theory into actionable optimization objectives. The target policy regularization penalizes deviation from theoretically sound strategies—using the known Nash equilibrium in simple games like Rock-Paper-Scissors, or the best-in-population policy in complex games like Leduc Poker. This provides direct gradient signals that guide policies toward equilibrium behavior rather than relying on opponent selection or meta-strategy computation.

\textbf{Exploitability-Driven Optimization:} The exploitability regularization term creates explicit pressure toward strategic robustness by penalizing policies that can be exploited by best-response opponents. This mechanism transforms population-based training data into direct optimization penalties, ensuring that robustness becomes an integral part of the learning objective rather than an emergent property of opponent diversity.

\textbf{Population-Based Training Architecture:} PRPO's population management system orchestrates coordinated training across multiple agents, implementing a two-phase regimen that balances internal robustness development through tournaments with external robustness through exploitative training. This approach systematically exposes policies to diverse strategic scenarios while maintaining computational efficiency through shared experience and parallel updates.

\textbf{Unified Framework Adaptability:} The framework's unified architecture enables seamless adaptation across game types through game-specific exploitability calculators, target policy generators, and exploiter opponent suites. This flexibility allows PRPO to maintain consistent methodology while adapting to the unique strategic characteristics of different game environments, from zero-sum normal-form games to imperfect-information extensive-form games.

\textbf{Convergence and Stability Improvements:} By embedding robustness objectives directly into the optimization process, PRPO achieves more stable convergence behavior compared to methods that rely solely on opponent sampling or meta-strategy computation. The explicit regularization terms provide consistent gradient signals that guide policies toward robust equilibria, reducing the sensitivity to training opponent selection and improving worst-case performance across diverse strategic scenarios.

\subsection{Trade-offs and Sensitivities}

While PRPO demonstrates significant improvements in strategic robustness, it introduces several trade-offs and sensitivities that warrant careful consideration.

\textbf{Computational Overhead and Oracle Costs:} The framework's strength in complex games comes at a computational price. Best-response oracle training requires substantial resources—particularly in imperfect-information games like Leduc Poker, where each oracle may need 250+ hands of training. This creates a clear trade-off between robustness gains and computational efficiency, with complex games requiring 2-3x more training time than standard approaches. The exploitability computation itself adds overhead, as fresh calculations must be performed before each policy update to maintain accurate gradient signals.

\textbf{Hyperparameter Sensitivity and Game Complexity:} PRPO exhibits varying sensitivity to its core hyperparameters across game types. In simple games like Rock-Paper-Scissors and Matching Pennies, the target regularization strength (\(\lambda_{\text{target}}\)) shows optimal performance in a broad range (1.0-2.0) with high robustness to variations. However, complex games like Leduc Poker require more conservative settings (\(\lambda_{\text{target}} = 0.2\)) to prevent gradient interference. The exploitability regularization strength (\(\lambda_{\text{exploit}}\)) demonstrates remarkable robustness across all tested ranges (0.5-2.0), suggesting the adaptive penalty mechanism successfully compensates for parameter variations.

\textbf{Architectural and Training Sensitivities:} The framework's performance can be sensitive to network initialization and exploration strategies. The custom Xavier initialization with small gain values (0.1) proves crucial for encouraging early exploration toward uniform policies, while higher entropy coefficients (0.05 vs. standard 0.01) maintain necessary stochasticity. Gradient clipping at 0.5 prevents instability from regularization terms, though this may occasionally limit the magnitude of game-theoretic corrections in challenging environments.

\textbf{Scalability and Game-Specific Adaptations:} While the unified architecture supports diverse game environments, performance scaling reveals interesting patterns. In zero-sum games with known analytical solutions, PRPO achieves near-optimal performance with minimal overhead. General-sum games like Stag Hunt benefit from the framework's ability to balance individual and social welfare objectives. However, imperfect-information games expose computational bottlenecks, where the cost of best-response computation may outweigh benefits for simpler architectures or training paradigms.

These trade-offs highlight that PRPO represents a sophisticated robustness optimization framework rather than a universally optimal solution. The framework's effectiveness depends on the specific game characteristics, available computational resources, and the researcher's tolerance for hyperparameter tuning. The key insight is that strategic robustness often requires explicit optimization, and PRPO provides a systematic approach to achieving this goal across diverse game environments.

\subsection{Limitations and Future Work}
Robustness gains in KP suggest architectural/training sensitivities; richer best-response oracles and longer horizons may further close gaps. Extending to multi-player and continuous-control domains is a natural next step.
\begin{itemize}
\item What Gauntlet reveals: evaluation matters; narrow metrics miss strategic brittleness.
\item Why PRPO helps: regularizers translate strategic desiderata into gradients.
\item Trade-offs: compute costs for oracles/BRs vs gains in robustness; sensitivity to \(\lambda\).
\item General-sum nuances: balancing social welfare with robustness; metric design implications.
\end{itemize}

%%%%%%%%%%%%%%%%%%%%%%%%%%%%%%%%%%%%%%%%%%%%%%%%%%%%%%%%%%%%%%%%%%%%%%%%

%%%%%%%%%%%%%%%%%%%%%%%%%%%%%%%%%%%%%%%%%%%%%%%%%%%%%%%%%%%%%%%%%%%%%%%%

\section{Conclusion}

This work makes fundamental contributions to the emerging field of robust multi-agent reinforcement learning by addressing a critical gap: the discrepancy between theoretical game-theoretic desiderata and practical policy optimization. We introduced two complementary innovations—the Gauntlet benchmark framework and Population-Regularized Policy Optimization (PRPO)—that together provide both rigorous evaluation tools and principled training methods for achieving strategic robustness.

\textbf{Summary of Contributions:} The Gauntlet benchmark represents a paradigm shift in MARL evaluation, moving beyond simplistic win-rate metrics to comprehensive robustness assessment. By systematically exposing policies to diverse challenger populations across multiple strategic dimensions, it reveals evaluation blind spots that standard approaches miss. Our empirical analysis demonstrates how conventional training paradigms can produce policies with high average performance but catastrophic failures against specific opponent types.

PRPO addresses these evaluation insights by embedding game-theoretic robustness directly into the optimization process. Unlike existing methods that rely on opponent selection or meta-strategy computation, PRPO transforms population-based training data into explicit regularization terms that guide policy optimization toward strategic equilibrium. The framework's unified architecture enables seamless adaptation across diverse game environments, from zero-sum normal-form games to complex imperfect-information scenarios.

\textbf{Empirical Validation:} Our comprehensive experiments across five diverse game environments provide strong evidence for the effectiveness of both contributions. PRPO achieves dramatic improvements in strategic robustness—reducing exploitability by 2-3 orders of magnitude in simple games like Rock-Paper-Scissors and Matching Pennies, while maintaining competitive performance in complex imperfect-information games like Leduc Poker. The Gauntlet framework successfully identifies these improvements, with composite robustness scores providing interpretable metrics that correlate with genuine strategic capabilities.

The ablation studies validate the design choices: target policy regularization provides consistent improvements across game types, while exploitability penalties offer complementary robustness benefits. Hyperparameter sensitivity analysis reveals that PRPO maintains effectiveness across a broad range of configurations, with conservative settings proving beneficial for complex games.

\textbf{Broader Implications:} This work establishes a new foundation for robust MARL research by demonstrating that strategic robustness requires both better evaluation frameworks and principled optimization methods. The Gauntlet benchmark provides a standardized evaluation protocol that can be adopted across the community, enabling fair comparison of robustness improvements and systematic identification of strategic weaknesses.

The success of PRPO suggests that explicit game-theoretic regularization represents a promising direction for robust learning beyond the specific methods presented here. By translating abstract strategic concepts into concrete optimization objectives, the framework opens new avenues for incorporating domain knowledge into reinforcement learning systems.

\textbf{Future Directions:} Several promising research directions emerge from this work. First, extending the Gauntlet framework to multi-agent scenarios beyond two-player games could provide insights into robustness in more complex social settings. Second, developing more sophisticated best-response oracles and meta-game computation methods could further enhance PRPO's performance in imperfect-information domains. Third, investigating the theoretical connections between explicit regularization and equilibrium convergence could provide deeper understanding of why these methods work.

The integration of PRPO with advanced neural architectures, such as transformers or graph neural networks, represents another promising direction for scaling robust MARL to real-world applications. Additionally, extending the framework to continuous action spaces and multi-objective optimization could broaden its applicability to domains like robotics, autonomous driving, and resource allocation.

\textbf{Legacy and Impact:} This work demonstrates that the path to robust multi-agent systems requires both better evaluation tools and principled training methods. By providing both, we hope to catalyze a shift in the field toward more rigorous evaluation practices and theoretically grounded optimization approaches. The open-source release of our framework ensures that these contributions can be built upon by the broader research community.

The fundamental insight—that strategic robustness requires explicit optimization rather than emergent properties of training dynamics—has implications extending beyond MARL to broader questions of reliable artificial intelligence. As AI systems increasingly operate in multi-agent and adversarial environments, the need for robust strategic reasoning becomes paramount. This work provides both conceptual foundations and practical tools for addressing this critical challenge. We encourage researchers to adopt these methods and build upon them to advance the frontier of robust multi-agent learning.

%%%%%%%%%%%%%%%%%%%%%%%%%%%%%%%%%%%%%%%%%%%%%%%%%%%%%%%%%%%%%%%%%%%%%%%%



%%% Balance columns on the final page for camera-ready
\balance

%%%%%%%%%%%%%%%%%%%%%%%%%%%%%%%%%%%%%%%%%%%%%%%%%%%%%%%%%%%%%%%%%%%%%%%%

\iffalse
\section{Introduction}

This document explains the main features of the `\texttt{aamas}' 
document class, which is essentially identical to the `\texttt{acmart}'
document class provided by the ACM. The only difference is a minor 
modification to allow for the correct copyright attribution to IFAAMAS.
For detailed documentation of the original document class, please refer
to the relevant website maintained by the~ACM:
%
\begin{center}
\url{https://www.acm.org/publications/proceedings-template}
\end{center}
%
The first command in your source file should be either one of these:
\begin{verbatim}
    \documentclass[sigconf,anonymous]{aamas}
    \documentclass[sigconf]{aamas}
\end{verbatim}
%
The first variant should be
used when you submit your paper for blind review; it will replace the names of the authors with the submission number.
The second variant should be used for final papers. 

Make sure your paper includes the correct copyright information and 
the correct specification of the \emph{ACM Reference Format}. Both of 
these will be generated automatically if you include the correct 
\emph{copyright block} as shown in the source file of this document.

Modifying the template---e.g., by changing margins, typeface sizes, 
line spacing, paragraph or list definitions---or making excessive use 
of the `\verb|\vspace|' command to manually adjust the vertical spacing 
between elements of your work is not allowed. You risk getting your 
submission rejected (or your final paper excluded from the proceedings) 
in case such modifications are discovered. The `\texttt{aamas}' document 
class requires the use of the \textit{Libertine} typeface family, which 
should be included with your \LaTeX\ installation. Please do not use 
other typefaces instead.

Please consult the \emph{Call for Papers} for information on matters 
such as the page limit or anonymity requirements. It is available from
the conference website:
%
\begin{center}
\url{https://cyprusconferences.org/aamas2026/}
\end{center}
%
To balance the columns on the final page of your paper, use the 
`\texttt{balance}' package and issue the `\verb|\balance|' command
 somewhere in the text of what would be the first column of the last 
 page without balanced columns. This will be required for final papers.

%%%%%%%%%%%%%%%%%%%%%%%%%%%%%%%%%%%%%%%%%%%%%%%%%%%%%%%%%%%%%%%%%%%%%%%%

\section{The Preamble}

You will be assigned a submission number when you register the abstract 
of your paper on \textit{OpenReview}. Include this number in your 
document using the `\verb|\acmSubmissionID|' command.

Then use the familiar commands to specify the title and authors of your
paper in the preamble of the document. The title should be appropriately 
capitalised (meaning that every `important' word in the title should 
start with a capital letter). For the final version of your paper, make 
sure to specify the affiliation and email address of each author using 
the appropriate commands. Specify an affiliation and email address 
separately for each author, even if two authors share the same 
affiliation. You can specify more than one affiliation for an author by 
using a separate `\verb|\affiliation|' command for each affiliation.

Provide a short abstract using the `\texttt{abstract}' environment.
 
Finally, specify a small number of keywords characterising your work, 
using the `\verb|\keywords|' command. 

%%%%%%%%%%%%%%%%%%%%%%%%%%%%%%%%%%%%%%%%%%%%%%%%%%%%%%%%%%%%%%%%%%%%%%%%

\section{The Body of the Paper}

For help with typesetting the body of your paper in \LaTeX\@, please 
make use of the familiar resources~\cite{Lam94}. In this section we 
merely highlight a few specific features. 

\subsection{Mathematical Expressions}

You can typeset all sorts of in-line mathematical expressions 
with the usual \verb|$...$| construct, as in 
$\Diamond\Diamond\varphi \rightarrow \Diamond\varphi$ or 
$\boldsymbol{R} = (R_1,\ldots,R_n)$.
For more complex expressions, it may often be preferable to use one of
the various equation-type environments available in \LaTeX\@, as shown 
in the following example:
%
\begin{eqnarray}
Z_i & = & \frac{u_i(x_i) - u_i(x_{-i})}{u_i(x_i)}
\end{eqnarray}
%
Here is a second example for an equation: 
%
\begin{eqnarray}\label{eq:vcg}
p_i(\boldsymbol{\hat{v}}) & = &
\sum_{j \neq i} \hat{v}_j(f(\boldsymbol{\hat{v}}_{-i})) - 
\sum_{j \neq i} \hat{v}_j(f(\boldsymbol{\hat{v}})) 
\end{eqnarray}
%
Use the usual combination of `\verb|\label|' and `\verb|\ref|' to refer
to numbered equations, such as Equation~(\ref{eq:vcg}) above. Of course,
introducing numbers in the first place is only helpful if you in fact 
need to refer back to the equation in question elsewhere in the paper.


\subsection{Tables and Figures}

Use the `\texttt{table}' environment (or its variant `\texttt{table*}')
in combination with the `\texttt{tabular}' environment to typeset tables
as floating objects. The `\texttt{aamas}' document class includes the 
`\texttt{booktabs}' package for preparing high-quality tables. Tables 
are often placed at the top of a page near their initial cite, as done 
here for Table~\ref{tab:locations}.

\begin{table}[t]
	\caption{Locations of the first five editions of AAMAS}
	\label{tab:locations}
	\begin{tabular}{rll}\toprule
		\textit{Year} & \textit{City} & \textit{Country} \\ \midrule
		2002 & Bologna & Italy \\
		2003 & Melbourne & Australia \\
		2004 & New York City & USA \\
		2005 & Utrecht & The Netherlands \\
		2006 & Hakodate & Japan \\ \bottomrule
	\end{tabular}
\end{table}

The caption of a table should be placed \emph{above} the table. 
Always use the `\verb|\midrule|' command to separate header rows from 
data rows, and use it only for this purpose. This enables assistive 
technologies to recognise table headers and support their users in 
navigating tables more easily.

%%% The following command should be issued somewhere in the first column 
%%% of the final page of your paper.
\balance

Use the `\texttt{figure}' environment for figures. If your figure 
contains third-party material, make sure to clearly identify it as such.
Every figure should include a caption, and this caption should be placed 
\emph{below} the figure itself, as shown here for Figure~\ref{fig:logo}.

\begin{figure}[h]
  \centering
  
\includegraphics[width=0.75\linewidth]{aamas2026logo}
  \caption{The logo of AAMAS 2026}
  \label{fig:logo}
  \Description{Logo of AAMAS 2026 -- The 25th International Conference on Autonomous Agents and Multiagent Systems.}
\end{figure}

In addition, every figure should also have a figure description, unless
it is purely decorative. Use the `\verb|\Description|' command for this 
purpose. These descriptions will not be printed but can be used to 
convey what's in an image to someone who cannot see it. They are also 
used by search engine crawlers for indexing images, and when images 
cannot be loaded. A figure description must consist of unformatted plain 
text of up to 2000~characters. For example, the definition of 
Figure~\ref{fig:logo} in the source file of this document includes the 
following description: ``Logo of AAMAS 2026 -- The 25th International Conference on Autonomous Agents and Multiagent Systems.'' For more information on how best to write figure descriptions 
and why doing so is important, consult the information available here: 
%
\begin{center}
\url{https://www.acm.org/publications/taps/describing-figures/}
\end{center}
%
The use of colour in figures and graphs is permitted, provided they 
remain readable when printed in greyscale and provided they are 
intelligible also for people with a colour vision deficiency.

%%%%%%%%%%%%%%%%%%%%%%%%%%%%%%%%%%%%%%%%%%%%%%%%%%%%%%%%%%%%%%%%%%%%%%%%

\section{Citations and References}
  
The use of the \BibTeX\ to prepare your list of references is highly 
recommended. To include the references at the end of your document, put 
the following two commands just before the `\verb|\end{document}|' 
command in your source file:
%
\begin{verbatim}
   \bibliographystyle{ACM-Reference-Format}
   \bibliography{sample}
\end{verbatim}
%
Here we assume that `\texttt{sample.bib}' is the name of your 
\BibTeX\ file. Use the `\verb|\cite|' command to produce citations 
to your references. Here are a few examples for citations of journal 
articles~\cite{GrKr96,WoJe95}, books~\cite{Knu97}, articles in 
conference proceedings~\cite{Hag1993}, technical reports~\cite{Har78},
Master's and PhD theses~\cite{Ani03,Cla85}, online videos~\cite{Oba08}, 
datasets~\cite{AnMC13}, and patents~\cite{Sci09}. Both citations and 
references are numbered by default. 

Make sure you provide complete and correct bibliographic information 
for all your references, and list authors with their full names 
(``Donald E.\ Knuth'') rather than just initials (``D.\ E.\ Knuth''). 

%%%%%%%%%%%%%%%%%%%%%%%%%%%%%%%%%%%%%%%%%%%%%%%%%%%%%%%%%%%%%%%%%%%%%%%%

%%% The acknowledgments section is defined using the "acks" environment
%%% (rather than an unnumbered section). The use of this environment 
%%% ensures the proper identification of the section in the article 
%%% metadata as well as the consistent spelling of the heading.

\fi
\begin{acks}
This work is prepared under anonymous submission. Detailed acknowledgments will be provided in the camera-ready version.
\end{acks}

%%%%%%%%%%%%%%%%%%%%%%%%%%%%%%%%%%%%%%%%%%%%%%%%%%%%%%%%%%%%%%%%%%%%%%%%

%%% The next two lines define, first, the bibliography style to be 
%%% applied, and, second, the bibliography file to be used.

\bibliographystyle{ACM-Reference-Format} 
\bibliography{sample}

%%%%%%%%%%%%%%%%%%%%%%%%%%%%%%%%%%%%%%%%%%%%%%%%%%%%%%%%%%%%%%%%%%%%%%%%

\end{document}

%%%%%%%%%%%%%%%%%%%%%%%%%%%%%%%%%%%%%%%%%%%%%%%%%%%%%%%%%%%%%%%%%%%%%%%%

